\chapter{KẾT LUẬN}
\label{Chapter6}

\section{Tóm tắt đóng góp}
% - Hệ thống end-to-end mã nguồn mở, dùng lại được cho các đề tài tiếp theo.
% - Kiến trúc phân tầng: firmware chỉ chuyển dữ liệu thô, server-side rule chain xử lý.
% - Bảo mật ở cả 2 lớp: HMAC-16 chống giả mạo beacon và MQTTS chống nghe lén.
% - Cơ chế OTA cho toàn hệ thống, không cần tháo lắp phần cứng.
% - Thuật toán hysteresis + leaky-bucket debounce chống flapping zone hiệu quả trong thí nghiệm.

\section{Hạn chế}
% - Chỉ định vị level phòng, không cho ra tọa độ (x,y).
% - Chỉ chạy trong 1 LAN, chưa xử lý đa site.
% - Chưa test scale > 10 tag đồng thời.
% - Chưa đo tiêu thụ pin dài hạn của tag.
% - Chưa test hiệu năng khi có nhiễu WiFi 2.4 GHz nặng.

\section{Hướng phát triển}

\subsection{Định vị tọa độ (x,y)}
% - Trilateration từ 3+ RSSI.
% - Kalman 2D hoặc particle filter.

\subsection{Học máy hiệu chỉnh path-loss}
% - Train mô hình theo môi trường cụ thể.
% - Adaptive n theo thời gian.

\subsection{Fleet management}
% - Quản lý hàng trăm gateway ở nhiều site.
% - Đồng bộ config qua cloud.

\subsection{Chống nhiễu 2.4 GHz}
% - Coexistence với WiFi.
% - Bluetooth 5 LE Coded PHY cho vùng phủ xa.
